% SPDX-License-Identifier: GPL-3.0-or-later
% Copyright (C) 2026 Alessandro Greco
\chapter{Порт управления и Stem}\label{cap:18-stem}
\aggiornatoal{23 сентября 2026 г.}{tor 0.4.9.13; control-spec; Stem 1.8.2}

\begin{obiettivi}
К концу этой главы вы сможете:
\begin{itemize}[nosep]
  \item взаимодействовать с tor через порт управления, в том числе вручную;
  \item использовать основные команды: \texttt{GETINFO}, \texttt{SETCONF}, \texttt{SIGNAL}, \file{ADD\_ONION} и подписку на события;
  \item писать программы на Python с библиотекой Stem;
  \item наблюдать изоляцию потоков и влияние новой идентичности на работающую систему (Лабораторная работа~11).
\end{itemize}
\end{obiettivi}

Tor Browser показывает цепь, используемую текущей страницей; такие инструменты, как Nyx, отображают пропускную способность и соединения реле; системы измерения сети строят цепи по запросу. Все эти приложения используют один интерфейс: \textbf{протокол управления} tor, доступный через \opt{ControlPort} (\cref{sec:torrc-controllo}).

% ---------------------------------------------------------------------------
\section{Протокол управления}\label{sec:stem-protocollo}

Протокол текстовый и напоминает почтовые протоколы~\spec{control-spec}{Message format}. Клиент отправляет по одной команде в строке, а tor отвечает строками, начинающимися с трёхзначного кода. Код \texttt{250} означает успех. Дефис после кода (\texttt{250-}) указывает, что ответ продолжается, а пробел (\texttt{250 }) — что это последняя строка; символ \texttt{+} вводит многострочный блок, завершаемый строкой, содержащей только точку. Коды \texttt{5xx} обозначают ошибки, а \texttt{650} — асинхронные события.

Попробуем это на клиенте лабораторной сети, открыв простое TCP-соединение с портом 8008:
\begin{lstlisting}[numbers=none]
> PROTOCOLINFO 1
250-PROTOCOLINFO 1
250-AUTH METHODS=COOKIE,SAFECOOKIE COOKIEFILE="/lab/net/.../008c/control_auth_cookie"
250-VERSION Tor="0.4.9.13"
250 OK
> AUTHENTICATE 52977550adebb73908610d7fa3355aba3fbdbdaef01d5ef86523d3232b87e5fa
250 OK
> GETINFO status/bootstrap-phase
250-status/bootstrap-phase=NOTICE BOOTSTRAP PROGRESS=100 TAG=done SUMMARY="Done"
250 OK
\end{lstlisting}
Это запись предыдущего демонстрационного запуска: пути, версия и cookie аутентификации зависят от конкретного экземпляра. \textbf{Никогда не публикуйте реальный cookie порта управления}: его обладатель может пройти аутентификацию.
\texttt{PROTOCOLINFO} входит в число немногих команд, разрешённых до аутентификации (остальные — \texttt{AUTHCHALLENGE}, \texttt{AUTHENTICATE} и \texttt{QUIT}). Она сообщает поддерживаемые методы аутентификации и путь к cookie-файлу. При аутентификации \texttt{COOKIE} клиент отправляет содержимое файла в шестнадцатеричном виде. \texttt{SAFECOOKIE} позволяет не передавать сам cookie: клиент и tor обмениваются случайными значениями и доказывают знание cookie с помощью HMAC, не позволяя программе, выдающей себя за tor, украсть его.

\begin{figure}[htbp]
\centering
% SPDX-License-Identifier: GPL-3.0-or-later
% Copyright (C) 2026 Alessandro Greco
%% Альтернативное описание для EPUB.
\altfig{Сокращённый пример сеанса протокола управления между приложением и tor. В AUTHENTICATE указан заполнитель для cookie в шестнадцатеричном виде, а не значение для копирования. После 250 OK на GETINFO circuit-status приходит многострочный ответ: заголовок 250+circuit-status=, строки данных, точка на отдельной строке и завершающее 250 OK. Команда SETEVENTS STREAM CIRC включает события, подписка подтверждается ответом 250 OK. Затем tor может асинхронно отправлять события 650 STREAM и 650 CIRC, в том числе между ответами на другие команды. Идентификаторы условные; PROTOCOLINFO и SAFECOOKIE не показаны.}
% Simplified control-port session with multiline GETINFO reply
\begin{tikzpicture}[font=\sffamily\scriptsize,>=Stealth,
  cmd/.style={->,torviola},
  ris/.style={->,torverde!60!black}]
  \node[client,minimum width=2.4cm,align=center] (a) at (0,0) {приложение\\(Stem, Nyx)};
  \node[client,minimum width=2.4cm,fill=torviola!20] (t) at (8.4,0) {tor};
  \draw[torgrigio] (0,-0.45) -- (0,-6.45); \draw[torgrigio] (8.4,-0.45) -- (8.4,-6.45);
  \draw[cmd] (0,-1.0) -- node[above,font=\ttfamily\tiny] {AUTHENTICATE [hex-cookie]} (8.4,-1.0);
  \draw[ris] (8.4,-1.6) -- node[above,font=\ttfamily\tiny] {250 OK} (0,-1.6);
  \draw[cmd] (0,-2.3) -- node[above,font=\ttfamily\tiny] {GETINFO circuit-status} (8.4,-2.3);
  \draw[ris] (8.4,-2.95) -- node[above,font=\ttfamily\tiny] {250+circuit-status=} (0,-2.95);
  \node[font=\ttfamily\tiny,align=center,text=torgrigio] at (4.2,-3.45) {7 BUILT ...\\.\\250 OK};
  \draw[cmd] (0,-4.25) -- node[above,font=\ttfamily\tiny] {SETEVENTS STREAM CIRC} (8.4,-4.25);
  \draw[ris] (8.4,-4.85) -- node[above,font=\ttfamily\tiny] {250 OK} (0,-4.85);
  \draw[ris] (8.4,-5.5) -- node[above,font=\ttfamily\tiny] {650 STREAM 42 SUCCEEDED ...} (0,-5.5);
  \draw[ris] (8.4,-6.15) -- node[above,font=\ttfamily\tiny] {650 CIRC 7 BUILT ...} (0,-6.15);
  \draw[decorate,decoration={brace,amplitude=3pt},torgrigio] (8.8,-5.35) -- (8.8,-6.3);
  \node[anchor=west,font=\sffamily\tiny,text=torgrigio,align=left] at (9.0,-5.85) {асинхронные события};
  \node[anchor=east,font=\sffamily\tiny,text=torgrigio] at (-0.3,-1.3) {аутентификация};
  \node[anchor=east,font=\sffamily\tiny,text=torgrigio] at (-0.3,-2.6) {запрос};
  \node[anchor=east,font=\sffamily\tiny,text=torgrigio] at (-0.3,-4.55) {подписка};
\end{tikzpicture}

\caption{Упрощённый сеанс протокола управления: после аутентификации COOKIE (в схеме указан заполнитель) запрос GETINFO circuit-status может вернуть многострочные данные, завершающиеся отдельной точкой и 250 OK. Подписка SETEVENTS подтверждается 250 OK; асинхронные события 650 могут перемежаться с ответами на другие команды. Идентификаторы условные.}
\label{fig:controllo}
\end{figure}

\subsection{Цепи с точки зрения tor}

\texttt{GETINFO circuit-status} перечисляет клиентские цепи. Сокращённая запись:
\begin{lstlisting}[numbers=none]
3 BUILT $50160F8F...~test004r,$12415FEB...~test005r,$9449F36E...~test007e
  BUILD_FLAGS=NEED_CAPACITY,NEED_UPTIME PURPOSE=CONFLUX_LINKED
  CONFLUX_ID=DE224696F58F72950D70DDA2D7981726 CONFLUX_RTT=42538
\end{lstlisting}
Каждое реле представлено отпечатком идентичности RSA и псевдонимом. Цепи Conflux появляются парами с одинаковым \file{CONFLUX\_ID}, причём две \enquote{ветви} начинаются с разных входных реле (guard). tor измеряет время прохождения туда и обратно (\file{CONFLUX\_RTT}, в микросекундах), чтобы решить, по какой ветви передавать данные (\cref{cap:31-futuro}). Одна деталь в коде tor (\file{src/feature/control/control\_fmt.c}) заслуживает внимания: \file{CONFLUX\_ID} раскрывает только \emph{половину} секретного идентификатора пары. Тот, кто знал бы идентификатор целиком, мог бы с позиции выходного реле добавлять ветви к паре и внедрять данные.

\subsection{Основные команды}

\begin{center}\small
\begin{tabularx}{\textwidth}{@{}l>{\raggedright\arraybackslash}X@{}}
\toprule
Команда & Назначение \\
\midrule
\texttt{GETINFO} & читает информацию: версию, состояние начальной загрузки, цепи, потоки, консенсус, дескрипторы\ldots \\
\texttt{GETCONF}, \texttt{SETCONF} & читает и изменяет текущую конфигурацию \\
\texttt{SIGNAL} & отправляет сигналы: \texttt{NEWNYM} (новая идентичность), \texttt{RELOAD}, \texttt{SHUTDOWN}\ldots \\
\texttt{SETEVENTS} & подписывает на события: \texttt{CIRC}, \texttt{STREAM}, \texttt{BW}, \file{HS\_DESC}\ldots \\
\texttt{EXTENDCIRCUIT} & строит цепь по маршруту, выбранному контроллером \\
\texttt{ATTACHSTREAM} & назначает поток цепи (вместе с \opt{\_\_LeaveStreamsUnattached}) \\
\file{ADD\_ONION}, \file{DEL\_ONION} & создаёт и удаляет onion-сервисы без редактирования \file{torrc} \\
\texttt{HSFETCH} & загружает дескриптор onion-сервиса (использовалась в \cref{sec:lab-onion}) \\
\bottomrule
\end{tabularx}
\end{center}

\begin{attenzione}[Контроль над портом означает контроль над tor]
Через порт управления можно указать tor, какие реле использовать, заставить его строить одношаговые цепи и раскрыть назначения всех потоков. Вредоносная программа с доступом к порту управления способна разрушить анонимность пользователя. Поэтому такие системы, как Tails и Whonix, помещают между приложениями и портом управления фильтр (onion-grater), разрешающий только строго необходимые команды.
\end{attenzione}

% ---------------------------------------------------------------------------
\section{Stem}\label{sec:stem-libreria}

\textbf{Stem} — Python-библиотека Tor Project для протокола управления. Она скрывает детали формата сообщений, разбирает ответы и события и предоставляет удобные объекты. Последняя версия, 1.8.2, вышла в июне 2023 года. Библиотека работает с актуальными версиями tor и используется многими инструментами, но новых релизов давно не было. Для новой функциональности иногда приходится отправлять команды напрямую (\texttt{Controller.msg}). У Arti другой интерфейс — Arti RPC (\cref{cap:30-arti}).

Минимальный пример, уже использованный в \cref{sec:rp-stem}:
\begin{lstlisting}[language=Python]
from stem.control import Controller

with Controller.from_port(port=8008) as c:
    c.authenticate()                   # найти cookie или запросить пароль
    print(c.get_version())
    for circ in c.get_circuits():
        print(circ.id, circ.purpose, [name for _, name in circ.path])
\end{lstlisting}

Для получения событий зарегистрируем функцию. Следующий пример --- \emph{самостоятельная программа}: слушатель регистрируется при открытом контроллере, после чего программа ждёт десять секунд для получения событий. Время ожидания здесь только демонстрационное.
\begin{lstlisting}[language=Python]
import time
from stem.control import Controller, EventType

def on_circuit(event):
    print(event.id, event.status, event.purpose)

with Controller.from_port(port=8008) as c:
    c.authenticate()
    c.add_event_listener(on_circuit, EventType.CIRC)
    time.sleep(10)
\end{lstlisting}
Stem вызывает эту функцию из своего потока обработки событий. Такой поток только один, и все зарегистрированные обработчики вызываются последовательно. Поэтому медленный обработчик событий задерживает последующие события. Если обработка требует значительных вычислений, передайте её отдельному потоку.

% ---------------------------------------------------------------------------
\section{Лабораторная работа: изоляция потоков в реальном времени}\label{sec:lab-stem}

\begin{laboratorio}[11]{Изоляция потоков и новая идентичность}
\textbf{Цель.} Использовать Stem на клиенте лабораторной сети, чтобы проверить, какие запросы совместно используют цепь в зависимости от SOCKS-учётных данных и что меняется после \texttt{SIGNAL NEWNYM}.\\
\textbf{Предварительные условия.} Лабораторная сеть должна быть запущена, а локальный веб-сервер — работать на порту 8080. Запустите сценарий \file{labs/ru/лаборатория/isolation.py} внутри контейнера, где доступны Stem и \cmd{curl}.\\
\textbf{Оценочное время.} 1 час.
\end{laboratorio}

\paragraph{Шаг 1: регистрация потоков.} Подпишемся на события \texttt{STREAM} и при каждом успешном потоке запишем цепь, к которой он был присоединён. Фильтрация по назначению обязательна: tor также открывает потоки для собственных задач, например загрузки документов каталога.
\begin{lstlisting}[language=Python]
assigned = []
def on_stream(ev):
    if ev.status == StreamStatus.SUCCEEDED and ev.target == DESTINATION:
        assigned.append((ev.id, ev.circ_id, ev.target))

c.add_event_listener(on_stream, EventType.STREAM)
\end{lstlisting}
Фрагмент взят из \file{labs/ru/лаборатория/isolation.py}: полный сценарий также задаёт \texttt{DESTINATION}, импортирует \texttt{StreamStatus} и держит контроллер \texttt{c} открытым. Отдельно этот фрагмент запускать нельзя.

\paragraph{Шаг 2: запросы с разными учётными данными.} Выполним четыре запроса через \cmd{curl}: два с SOCKS-учётными данными \texttt{site-a.example:x1}, один с \texttt{site-b.example:x2} и один без учётных данных. Tor Browser использует домен основного сайта как имя пользователя SOCKS для каждой вкладки:
\begin{lstlisting}[language=bash,numbers=none]
$ curl --proxy-user site-a.example:x1 --socks5-hostname 127.0.0.1:9008 http://127.0.0.1:8080/
\end{lstlisting}

\paragraph{Шаг 3: результаты.} Ниже приведён вывод предыдущего эталонного запуска; идентификаторы Conflux, маршруты и задержка NEWNYM не являются постоянными значениями.
\begin{lstlisting}[numbers=none]
учётные данные site-a.example  -> conflux AADCE71E  ['test003r', 'test004r', 'test007e']
учётные данные site-a.example  -> conflux AADCE71E  ['test003r', 'test004r', 'test007e']
учётные данные site-b.example  -> conflux 31B1E3E0  ['test003r', 'test004r', 'test007e']
учётные данные -               -> conflux 35FD16D0  ['test003r', 'test004r', 'test007e']
(NEWNYM задержан на 6 с)
после NEWNYM, site-a.example -> conflux F7011ABE
\end{lstlisting}
Два запроса к \enquote{сайту A} используют один и тот же туннель, хотя выполняются в разное время. Запрос к \enquote{сайту B} и запрос без аутентификации используют другие туннели. Выведенный маршрут — одна из двух ветвей туннеля; в данном запуске они совпадают, поскольку в лабораторной сети мало реле. Далее объясняется, почему следует анализировать туннель, а не отдельную цепь.

\begin{approfondimento}[Одной цепи недостаточно: туннели Conflux]
Если повторить лабораторную работу много раз, два запроса с одинаковыми учётными данными иногда оказываются на цепях с разными номерами. Это не нарушение изоляции. Запрос полного списка цепей tor показывает, что почти все клиентские цепи имеют назначение \file{CONFLUX\_LINKED} и появляются \emph{парами} с одинаковым \file{CONFLUX\_ID}. Это две \enquote{ветви} одного туннеля (\cref{sec:fut-conflux}), и событие \texttt{STREAM} может указать любую из них. Следовательно, изоляция разделяет \emph{туннели}, а не отдельные ветви. Когда это возможно, сценарий книги группирует их по \file{CONFLUX\_ID}.

Нужно соблюдать ещё две меры предосторожности, обе можно проверить по журналам:
\begin{itemize}
  \item \texttt{SIGNAL NEWNYM} исполняется не чаще одного раза в десять секунд (\file{MAX\_SIGNEWNYM\_RATE} в \file{src/core/mainloop/mainloop.c}). Более ранний запрос откладывается, и tor записывает \enquote{Rate limiting NEWNYM request: delaying by N second(s)}. Если не дождаться выполнения, NEWNYM из предыдущего эксперимента может сработать уже в середине следующего.
  \item События \texttt{STREAM} также сообщают о потоках, которые tor открывает для собственных нужд, например для загрузки документов каталога. Фильтруйте их по назначению.
\end{itemize}
Даже при соблюдении этих мер в небольшой сети под высокой нагрузкой tor может выбрать только что построенный туннель для второго запроса. Поэтому автоматическая проверка книги повторяет эксперимент.
\end{approfondimento}

\begin{attenzione}[Ошибка в лабораторной работе]
В первом запуске все четыре запроса использовали одну цепь, что, казалось, показывало отсутствие изоляции. Проблема была в нашем сценарии: мы поместили \texttt{--proxy-user} сразу после опции \texttt{-o} команды \cmd{curl}, и программа восприняла его как имя выходного файла. Учётные данные вообще не отправлялись. Прежде чем делать вывод, что система ведёт себя не так, как предписывает спецификация, убедитесь, что эксперимент действительно делает то, что вы предполагаете.
\end{attenzione}

% ---------------------------------------------------------------------------
\section*{Итоги}
\begin{itemize}
  \item Протокол управления — текстовый: однострочные команды, трёхзначные коды ответов и асинхронные события с кодом 650.
  \item Используйте файл cookie, предпочтительно с методом \texttt{SAFECOOKIE}, или пароль. Тот, кто контролирует порт, контролирует tor.
  \item Stem предоставляет Python-интерфейс к командам и событиям. Библиотека стабильна, но новых релизов не было с 2023 года.
  \item Изоляция по SOCKS-учётным данным выделяет каждому сайту отдельный туннель, который при Conflux состоит из двух цепей. \texttt{NEWNYM} переключает на новые туннели не чаще одного раза в десять секунд.
\end{itemize}

% ---------------------------------------------------------------------------
\section*{Упражнения}

\begin{esercizio}[1]\label{es:18-codici}
Интерпретируйте ответ \texttt{552 Unrecognized key "pippo"} на команду \texttt{GETINFO pippo}. Чем код \texttt{5xx} отличается от события \texttt{650}?
\end{esercizio}

\begin{esercizio}[2]\label{es:18-onion}
С помощью Stem создайте временный onion-сервис функцией \file{create\_ephemeral\_hidden\_service} на клиенте лабораторной сети, перенаправив порт 80 на локальный сервер порта 8080. Подключитесь к нему через \cmd{curl}. Где хранятся его ключи?
\end{esercizio}

\begin{esercizio}[2]\label{es:18-banda}
Напишите программу, которая каждую секунду отображает входящую и исходящую пропускную способность клиента по событиям \texttt{BW}, пока через лабораторную сеть загружается большой файл.
\end{esercizio}

\begin{esercizio}[3]\label{es:18-percorso}
Используя \opt{\_\_LeaveStreamsUnattached} 1, \texttt{EXTENDCIRCUIT} и \texttt{ATTACHSTREAM}, направьте запрос по выбранному вами маршруту, например через \texttt{test004r}, \texttt{test005r}, \texttt{test006e}. Почему такая возможность опасна для анонимности при реальном использовании?
\end{esercizio}

% ---------------------------------------------------------------------------
\section*{Дополнительная литература}
\begin{itemize}
  \item \titolo{Tor's control protocol} (control-spec)~\cite{torspec}: все команды и события.
  \item Документация Stem: \url{https://stem.torproject.org/}.
\end{itemize}
